%TCIDATA{LaTeXparent=0,0,EOMN.tex}
                      

\subsection{General Form of Partitioning}

Consider the matrix equation 
\begin{equation}
\left( 
\begin{array}{cc}
{\Bbb A}_{11} & {\Bbb A}_{12} \\ 
{\Bbb A}_{21} & {\Bbb A}_{22}%
\end{array}%
\right) \left( 
\begin{array}{c}
{\bf x}_{1} \\ 
{\bf x}_{2}%
\end{array}%
\right) =\left( 
\begin{array}{c}
{\bf y}_{1} \\ 
{\bf y}_{2}%
\end{array}%
\right)  \label{ME1}
\end{equation}%
that is equivalent to the coupled equations 
\begin{eqnarray}
{\Bbb A}_{11}{\bf x}_{1}+{\Bbb A}_{12}{\bf x}_{2} &=&{\bf y}_{1}  \nonumber
\\
{\Bbb A}_{21}{\bf x}_{1}+{\Bbb A}_{22}{\bf x}_{2} &=&{\bf y}_{2}
\end{eqnarray}%
we can solve for ${\bf x}_{1}$ from the top equation to obtain a ''fixed'' $%
{\bf x}_{1}$ 
\begin{equation}
{\bf x}_{1}={\Bbb A}_{11}^{-1}\left( {\bf y}_{1}-{\Bbb A}_{12}{\bf x}%
_{2}\right)  \label{x1eq}
\end{equation}%
and then substitute in the bottom equation to obtain a solution for ${\bf x}%
_{2}$ 
\begin{equation}
{\bf x}_{2}=\left( {\Bbb A}_{22}+{\Bbb A}_{21}{\Bbb A}_{11}^{-1}{\Bbb A}%
_{12}\right) ^{-1}\left( {\bf y}_{2}-{\Bbb A}_{21}{\Bbb A}_{11}^{-1}{\bf y}%
_{1}\right)  \label{x2sol}
\end{equation}%
which then determines ${\bf x}_{1}$ by eq. $\left( \ref{x1eq}\right) $. The
choice of which variables to call ${\bf x}_{1}$ and ${\bf x}_{2}$ in any
problem is of course quite arbitrary. One should also observe that the ${\bf %
x}_{2}$ obtained from eq. $\left( \ref{x2sol}\right) $ can also be expressed
as 
\begin{equation}
{\bf x}_{2}={\Bbb A}_{22}^{-1}\left( {\bf y}_{2}-{\Bbb A}_{21}{\bf x}%
_{1}\right)  \label{x2leq}
\end{equation}%
where ${\bf x}_{1}$ is given by eq. $\left( \ref{x1eq}\right) $. If neither $%
{\Bbb A}$ nor ${\bf y}$ depend on ${\bf x}$ and all the inverses exist then
a fixed solution for ${\bf x}_{1}$ and ${\bf x}_{2}$ is obtained in one
iteration.

The equations we deal with are, however, in general nonlinear i.e. 
\begin{eqnarray}
{\Bbb A}_{ij} &=&{\Bbb A}_{ij}\left( {\bf x}_{1},{\bf x}_{2}\right) ;\quad
1\leq i,j\leq 2  \nonumber \\
{\bf y}_{i} &=&{\bf y}_{i}\left( {\bf x}_{1},{\bf x}_{2}\right) ;\quad 1\leq
i\leq 2
\end{eqnarray}%
and we have to iterate to obtain a solution by fixing a value for ${\bf x}%
_{1}$ and solving for ${\bf x}_{2\text{ }}$to obtain 
\begin{equation}
{\bf x}_{2}=\left( {\Bbb A}_{22}+{\Bbb A}_{21}{\Bbb A}_{11}^{-1}{\Bbb A}%
_{12}\right) ^{-1}\left( {\bf y}_{2}-{\Bbb A}_{21}{\Bbb A}_{11}^{-1}{\bf y}%
_{1}\right) ={\Bbb A}_{22}^{-1}\left( {\bf y}_{2}-{\Bbb A}_{21}{\bf x}%
_{1}\right)  \label{x2parteq}
\end{equation}%
until one reaches self consistency or fixing a value for ${\bf x}_{2}$ and
solving for ${\bf x}_{1\text{ }}$to obtain%
\begin{equation}
{\bf x}_{1}={\Bbb A}_{11}^{-1}\left( {\bf y}_{1}-{\Bbb A}_{12}{\bf x}%
_{2}\right)  \label{x1parteq}
\end{equation}%
until one reaches self consistency. When we have a partial self consistent
solution for ${\bf x}_{1}$ and use a fixed ${\bf x}_{2}$ we can substitute
back into the equations to obtain 
\begin{eqnarray}
&&\left( 
\begin{array}{cc}
{\Bbb A}_{11} & {\Bbb A}_{12} \\ 
{\Bbb A}_{21} & {\Bbb A}_{22}%
\end{array}%
\right) \left( 
\begin{array}{c}
{\Bbb A}_{11}^{-1}\left( {\bf y}_{1}-{\Bbb A}_{12}{\bf x}_{2}\right) \\ 
{\bf x}_{2}%
\end{array}%
\right) -\left( 
\begin{array}{c}
{\bf y}_{1} \\ 
{\bf y}_{2}%
\end{array}%
\right)  \nonumber \\
&=&\left( 
\begin{array}{c}
{\Bbb A}_{11}{\Bbb A}_{11}^{-1}\left( {\bf y}_{1}-{\Bbb A}_{12}{\bf x}%
_{2}\right) +{\Bbb A}_{12}{\bf x}_{2} \\ 
{\Bbb A}_{21}{\Bbb A}_{11}^{-1}\left( {\bf y}_{1}-{\Bbb A}_{12}{\bf x}%
_{2}\right) +{\Bbb A}_{22}{\bf x}_{2}%
\end{array}%
\right) -\left( 
\begin{array}{c}
{\bf y}_{1} \\ 
{\bf y}_{2}%
\end{array}%
\right)  \nonumber \\
&=&\left( 
\begin{array}{c}
\left( {\bf y}_{1}-{\Bbb A}_{12}{\bf x}_{2}\right) +{\Bbb A}_{12}{\bf x}_{2}
\\ 
{\Bbb A}_{21}{\Bbb A}_{11}^{-1}\left( {\bf y}_{1}-{\Bbb A}_{12}{\bf x}%
_{2}\right) +{\Bbb A}_{22}{\bf x}_{2}%
\end{array}%
\right) -\left( 
\begin{array}{c}
{\bf y}_{1} \\ 
{\bf y}_{2}%
\end{array}%
\right)  \nonumber \\
&=&\left( 
\begin{array}{c}
{\bf 0} \\ 
{\Bbb A}_{21}{\Bbb A}_{11}^{-1}\left( {\bf y}_{1}-{\Bbb A}_{12}{\bf x}%
_{2}\right) +{\Bbb A}_{22}{\bf x}_{2}-{\bf y}_{2}%
\end{array}%
\right)  \label{baksub1}
\end{eqnarray}%
While when we have a partial self consistent solution for ${\bf x}_{2}$ and
use a fixed ${\bf x}_{1}$ we can substitute back into the equations to
obtain 
\begin{eqnarray}
&&\left( 
\begin{array}{cc}
{\Bbb A}_{11} & {\Bbb A}_{12} \\ 
{\Bbb A}_{21} & {\Bbb A}_{22}%
\end{array}%
\right) \left( 
\begin{array}{c}
{\bf x}_{1} \\ 
\left( {\Bbb A}_{22}+{\Bbb A}_{21}{\Bbb A}_{11}^{-1}{\Bbb A}_{12}\right)
^{-1}\left( {\bf y}_{2}-{\Bbb A}_{21}{\Bbb A}_{11}^{-1}{\bf y}_{1}\right)%
\end{array}%
\right) -\left( 
\begin{array}{c}
{\bf y}_{1} \\ 
{\bf y}_{2}%
\end{array}%
\right)  \nonumber \\
&=&\left( 
\begin{array}{c}
{\Bbb A}_{11}{\bf x}_{1}+{\Bbb A}_{12}\left( {\Bbb A}_{22}+{\Bbb A}_{21}%
{\Bbb A}_{11}^{-1}{\Bbb A}_{12}\right) ^{-1}\left( {\bf y}_{2}-{\Bbb A}_{21}%
{\Bbb A}_{11}^{-1}{\bf y}_{1}\right) \\ 
{\Bbb A}_{21}{\bf x}_{1}+{\Bbb A}_{22}\left( {\Bbb A}_{22}+{\Bbb A}_{21}%
{\Bbb A}_{11}^{-1}{\Bbb A}_{12}\right) ^{-1}\left( {\bf y}_{2}-{\Bbb A}_{21}%
{\Bbb A}_{11}^{-1}{\bf y}_{1}\right)%
\end{array}%
\right) -\left( 
\begin{array}{c}
{\bf y}_{1} \\ 
{\bf y}_{2}%
\end{array}%
\right)  \nonumber \\
&=&\left( 
\begin{array}{c}
{\Bbb A}_{11}{\bf x}_{1}+{\Bbb A}_{12}{\Bbb A}_{22}^{-1}\left\{ {\bf y}_{2}-%
{\Bbb A}_{21}{\bf x}_{1}\right\} \\ 
{\bf y}_{2}%
\end{array}%
\right) -\left( 
\begin{array}{c}
{\bf y}_{1} \\ 
{\bf y}_{2}%
\end{array}%
\right)  \nonumber \\
&=&\left( 
\begin{array}{c}
{\Bbb A}_{11}{\bf x}_{1}+{\Bbb A}_{12}{\Bbb A}_{22}^{-1}\left\{ {\bf y}_{2}-%
{\Bbb A}_{21}{\bf x}_{1}\right\} -{\bf y}_{1} \\ 
{\bf 0}%
\end{array}%
\right)  \label{baksub2}
\end{eqnarray}%
The partitioning displayed above obviously only works when the various
inverses exist.

We can obtain equivalent equations by multiplying both sides of eq. $\left( %
\ref{ME1}\right) $ by an invertible matrix e.g. 
\begin{equation}
\left( 
\begin{array}{cc}
{\Bbb B}_{11} & {\Bbb B}_{12} \\ 
{\Bbb B}_{21} & {\Bbb B}_{22}%
\end{array}%
\right) \left( 
\begin{array}{cc}
{\Bbb A}_{11} & {\Bbb A}_{12} \\ 
{\Bbb A}_{21} & {\Bbb A}_{22}%
\end{array}%
\right) \left( 
\begin{array}{c}
{\bf x}_{1} \\ 
{\bf x}_{2}%
\end{array}%
\right) =\left( 
\begin{array}{cc}
{\Bbb B}_{11} & {\Bbb B}_{12} \\ 
{\Bbb B}_{21} & {\Bbb B}_{22}%
\end{array}%
\right) \left( 
\begin{array}{c}
{\bf y}_{1} \\ 
{\bf y}_{2}%
\end{array}%
\right) =\left( 
\begin{array}{c}
{\bf z}_{1} \\ 
{\bf z}_{2}%
\end{array}%
\right)  \label{pre}
\end{equation}%
and then partitioning in the same manner. If ${\Bbb A}_{12}={\Bbb A}_{21}=0$
then the equations are uncoupled if they are linear, but are still coupled
if they are nonlinear. If ${\Bbb A}_{11}={\Bbb A}_{22}=0$ then one can
premultiply eq. $\left( \ref{pre}\right) $ by $\left( 
\begin{array}{cc}
{\Bbb O} & {\Bbb I} \\ 
{\Bbb I} & {\Bbb O}%
\end{array}%
\right) $ and then solve, but then the partial solutions do not give the
same back substituted properties displayed in eqs. $\left( \ref{baksub1}%
\right) $ and $\left( \ref{baksub2}\right) $ i.e. the residual vectors of
the original problem generally have zero components in different positions
or no zero components at all.

\subsection{Zero ${\Bbb A}_{22}$ Block}

The {\em nonlinear }equation eq. $\left( \ref{ME1}\right) $ for ${\bf x}_{1}$
and ${\bf x}_{2}$ 
\begin{equation}
\left( 
\begin{array}{cc}
{\Bbb A}_{11} & {\Bbb A}_{12} \\ 
{\Bbb A}_{21} & {\Bbb A}_{22}%
\end{array}%
\right) \left( 
\begin{array}{c}
{\bf x}_{1} \\ 
{\bf x}_{2}%
\end{array}%
\right) -\left( 
\begin{array}{c}
{\bf y}_{1} \\ 
{\bf y}_{2}%
\end{array}%
\right) =0
\end{equation}%
does not have be solved in the preceding manner, in fact when one of the
diagonal blocks is zero it may become impossible. In particular consider 
\begin{equation}
\left( 
\begin{array}{cc}
0 & {\Bbb A}_{12} \\ 
{\Bbb A}_{21} & {\Bbb A}_{22}%
\end{array}%
\right) \left( 
\begin{array}{c}
{\bf x}_{1} \\ 
{\bf x}_{2}%
\end{array}%
\right) -\left( 
\begin{array}{c}
{\bf y}_{1} \\ 
{\bf y}_{2}%
\end{array}%
\right) =0\Longleftrightarrow \left. 
\begin{array}{c}
\qquad \qquad {\Bbb A}_{12}{\bf x}_{2}-{\bf y}_{1}=0 \\ 
{\Bbb A}_{21}{\bf x}_{1}+{\Bbb A}_{22}{\bf x}_{2}-{\bf y}_{2}=0%
\end{array}%
\right\}
\end{equation}%
which gives for fixed ${\bf x}_{1}$ the equation 
\begin{equation}
{\Bbb A}_{22}{\bf x}_{2}={\bf y}_{2}-{\Bbb A}_{21}{\bf x}_{1}  \label{x2le}
\end{equation}%
for ${\bf x}_{2}$. A self consistent solution, ${\bf x}_{2\#}\left( {\bf x}%
_{1\left( k\right) }\right) $, is obtained by iterating ${\bf x}_{2}$ in eq. 
$\left( \ref{x2le}\right) $ at the fixed value ${\bf x}_{1}={\bf x}_{1\left(
k\right) }$ 
\begin{equation}
{\bf x}_{2\#}\left( {\bf x}_{1\left( k\right) }\right) ={\Bbb A}%
_{22}^{-1}\left( {\bf y}_{2}-{\Bbb A}_{21}{\bf x}_{1\left( k\right) }\right)
\label{x2}
\end{equation}%
If one substitutes back this partial solution, noting that all quantities, $%
{\Bbb A}_{12},{\Bbb A}_{21},{\Bbb A}_{22},{\bf y}_{1}$ and ${\bf y}_{2}$ are
evaluated at $\left( {\bf x}_{1\left( k\right) },{\bf x}_{2\#}\left( {\bf x}%
_{1\left( k\right) }\right) \right) $, we obtain 
\begin{eqnarray}
&&\left( 
\begin{array}{cc}
0 & {\Bbb A}_{12} \\ 
{\Bbb A}_{21} & {\Bbb A}_{22}%
\end{array}%
\right) \left( 
\begin{array}{c}
{\bf x}_{1\left( k\right) } \\ 
{\Bbb A}_{22}^{-1}\left( {\bf y}_{2}-{\Bbb A}_{21}{\bf x}_{1\left( k\right)
}\right)%
\end{array}%
\right) -\left( 
\begin{array}{c}
{\bf y}_{1} \\ 
{\bf y}_{2}%
\end{array}%
\right) =  \nonumber \\
&&\qquad \left( 
\begin{array}{c}
{\Bbb A}_{12}{\Bbb A}_{22}^{-1}\left( {\bf y}_{2}-{\Bbb A}_{21}{\bf x}%
_{1\left( k\right) }\right) -{\bf y}_{1} \\ 
{\Bbb A}_{21}{\bf x}_{1\left( k\right) }+{\bf y}_{2}-{\Bbb A}_{21}{\bf x}%
_{1\left( k\right) }-{\bf y}_{2}%
\end{array}%
\right) =\left( 
\begin{array}{c}
{\Bbb A}_{12}{\Bbb A}_{22}^{-1}\left( {\bf y}_{2}-{\Bbb A}_{21}{\bf x}%
_{1\left( k\right) }\right) -{\bf y}_{1} \\ 
{\bf 0}%
\end{array}%
\right)  \label{baksub}
\end{eqnarray}%
The self consistent solution, ${\bf x}_{1\#\left( k+1\right) }$, for ${\bf x}%
_{1}$ at ${\bf x}_{2\#}\left( {\bf x}_{1\left( k\right) }\right) $ is then
obtained by iterating ${\bf x}_{1\left( k+1\right) }$ in%
\begin{equation}
{\Bbb A}_{12}{\Bbb A}_{22}^{-1}\left( {\bf y}_{2}-{\Bbb A}_{21}{\bf x}%
_{1\left( k+1\right) }\right) -{\bf y}_{1}=0
\end{equation}%
or in the formal relationship 
\begin{equation}
{\bf x}_{1\left( k+1\right) }={\Bbb A}_{12}^{-1}{\Bbb A}_{22}{\Bbb A}%
_{21}^{-1}\left( {\bf y}_{1}-{\Bbb A}_{12}{\Bbb A}_{22}^{-1}{\bf y}%
_{2}\right)  \label{x1kplus1}
\end{equation}%
where all quantities are evaluated at $\left( {\bf x}_{1\left( k+1\right) },%
{\bf x}_{2\#}\left( {\bf x}_{1\left( k\right) }\right) \right) $ and left
inverses are used for ${\Bbb A}_{12}$ and ${\Bbb A}_{21}$. These two coupled
iterative sequences can then be repeated until self consistency is reached
for both ${\bf x}_{1}$ and ${\bf x}_{2}$ to obtain a solution of the full
non linear problem.

If ${\Bbb A}_{22}$ is singular then eq. $\left( \ref{x2}\right) $ becomes%
\begin{equation}
{\bf x}_{2\#}\left( {\bf x}_{1\left( k\right) },{\bf k}\right) ={\Bbb A}%
_{22}^{-1}\left( {\bf y}_{2}-{\Bbb A}_{21}{\bf x}_{1\left( k\right) }\right)
+{\bf k}  \label{x2eq}
\end{equation}%
where we use the generalized inverse ${\Bbb A}_{22}^{-1}$ of ${\Bbb A}_{22}$%
, ${\bf k}\in \ker {\Bbb A}_{22}$ and the subsidiary condition 
\begin{equation}
P_{\ker {\Bbb A}_{22}}\left( {\bf y}_{2}-{\Bbb A}_{21}{\bf x}_{1\left(
k\right) }\right) =0  \label{x2cond}
\end{equation}%
must be satisfied, with all terms evaluated at $\left( {\bf x}_{1\left(
k\right) },{\bf x}_{2\#}\left( {\bf x}_{1\left( k\right) },{\bf k}\right)
\right) $. Then the back substitution eq. $\left( \ref{baksub}\right) $
becomes 
\begin{eqnarray}
&&\left( 
\begin{array}{cc}
0 & {\Bbb A}_{12} \\ 
{\Bbb A}_{21} & {\Bbb A}_{22}%
\end{array}%
\right) \left( 
\begin{array}{c}
{\bf x}_{1\left( k\right) } \\ 
{\Bbb A}_{22}^{-1}\left( {\bf y}_{2}-{\Bbb A}_{21}{\bf x}_{1\left( k\right)
}\right) +{\bf k}%
\end{array}%
\right) -\left( 
\begin{array}{c}
{\bf y}_{1} \\ 
{\bf y}_{2}%
\end{array}%
\right) =  \nonumber \\
&&\qquad \left( 
\begin{array}{c}
{\Bbb A}_{12}{\Bbb A}_{22}^{-1}\left( {\bf y}_{2}-{\Bbb A}_{21}{\bf x}%
_{1\left( k\right) }\right) +{\Bbb A}_{12}{\bf k}-{\bf y}_{1} \\ 
P_{\ker {\Bbb A}_{22}}^{c}\left( {\bf y}_{2}-{\Bbb A}_{21}{\bf x}_{1\left(
k\right) }\right) +{\Bbb A}_{21}{\bf x}_{1\left( k\right) }-{\bf y}_{2}%
\end{array}%
\right) =  \nonumber \\
&&\qquad \qquad \qquad \qquad \left( 
\begin{array}{c}
{\Bbb A}_{12}{\Bbb A}_{22}^{-1}\left( {\bf y}_{2}-{\Bbb A}_{21}{\bf x}%
_{1\left( k\right) }\right) +{\Bbb A}_{12}{\bf k}-{\bf y}_{1} \\ 
{\bf 0}%
\end{array}%
\right)  \label{subbak}
\end{eqnarray}%
for any ${\bf k}$. The equation for a self consistent value, ${\bf x}%
_{1\#\left( k+1\right) }\left( {\bf x}_{2\#}\left( {\bf x}_{1\left( k\right)
},{\bf k}\right) \right) $, for ${\bf x}_{1}$ at ${\bf x}_{2}={\bf x}%
_{2\#}\left( {\bf x}_{1\left( k\right) },{\bf k}\right) $ obtained from eq. $%
\left( \ref{subbak}\right) $ is%
\begin{equation}
{\Bbb A}_{12}{\Bbb A}_{22}^{-1}{\Bbb A}_{21}{\bf x}_{1\#\left( k+1\right) }=%
{\Bbb A}_{12}{\Bbb A}_{22}^{-1}{\bf y}_{2}-{\bf y}_{1}+{\Bbb A}_{12}{\bf k}
\label{scx1}
\end{equation}%
which can be formally solved as 
\begin{equation}
{\bf x}_{1\#\left( k+1\right) }\left( {\bf x}_{1\left( k\right) },{\bf k}%
\right) ={\Bbb A}_{21}^{-1}{\Bbb A}_{22}{\Bbb A}_{12}^{-1}\left( {\Bbb A}%
_{12}{\Bbb A}_{22}^{-1}{\bf y}_{2}-{\bf y}_{1}\right) ={\Bbb A}_{21}^{-1}%
{\bf y}_{2}-{\Bbb A}_{21}^{-1}{\Bbb A}_{22}{\Bbb A}_{12}^{-1}{\bf y}_{1}
\end{equation}%
where ${\Bbb A}_{12},{\Bbb A}_{21},{\Bbb A}_{22},{\bf y}_{1},$ and ${\bf y}%
_{2}$ are evaluated at $\left( {\bf x}_{1\#\left( k+1\right) },{\bf x}%
_{2\#}\left( {\bf x}_{1\left( k\right) },{\bf k}\right) \right) $. This
value can then be substituted back into the equation for ${\bf x}_{2}$,
which can then be iterated to self consistency once more, the new self
consistent value of ${\bf x}_{2}$ can then be used to obtain a new self
consistent value of ${\bf x}_{1}$ and so on.

Instead\ of iterating to self consistency for a new\ ${\bf x}_{1\#\left(
k+1\right) }$ and then substituting back into the equation for ${\bf x}_{2}$%
, which would be the most efficient way to solve the total coupled problem,
one could find paths of self consistent solutions, ${\bf x}_{2\#}\left( {\bf %
x}_{1},{\bf k}\right) $, for all possible ${\bf x}_{1}$'s producing paths, $%
{\bf x}_{2\#}\left( {\bf x}_{1},{\bf k}\right) $, of ${\bf x}_{2}$'s that
satisfy 
\begin{equation}
{\bf x}_{2\#}\left( {\bf x}_{1},{\bf k}\right) ={\Bbb A}_{22}^{-1}\left( 
{\bf y}_{2}-{\Bbb A}_{21}{\bf x}_{1}\right) +{\bf k}
\end{equation}%
where 
\begin{eqnarray}
{\Bbb A}_{22} &=&{\Bbb A}_{22}\left( {\bf x}_{1},{\bf x}_{2\#}\left( {\bf x}%
_{1},{\bf k}\right) \right)  \nonumber \\
{\bf y}_{2} &=&{\bf y}_{2}\left( {\bf x}_{1},{\bf x}_{2\#}\left( {\bf x}_{1},%
{\bf k}\right) \right)  \nonumber \\
{\Bbb A}_{21} &=&{\Bbb A}_{21}\left( {\bf x}_{1},{\bf x}_{2\#}\left( {\bf x}%
_{1},{\bf k}\right) \right)  \nonumber \\
{\bf k} &=&{\bf k}\left( {\bf x}_{1}\right)
\end{eqnarray}%
and 
\begin{eqnarray}
\left( 
\begin{array}{cc}
0 & {\Bbb A}_{12} \\ 
{\Bbb A}_{21} & {\Bbb A}_{22}%
\end{array}%
\right) \left( 
\begin{array}{c}
{\bf x}_{1} \\ 
{\bf x}_{2\#}\left( {\bf x}_{1},{\bf k}\right)%
\end{array}%
\right) -\left( 
\begin{array}{c}
{\bf y}_{1} \\ 
{\bf y}_{2}%
\end{array}%
\right) &=&\left( 
\begin{array}{c}
{\Bbb A}_{12}{\bf x}_{2\#}\left( {\bf x}_{1},{\bf k}\right) -{\bf y}_{1} \\ 
{\Bbb A}_{21}{\bf x}_{1}+{\Bbb A}_{22}{\bf x}_{2\#}\left( {\bf x}_{1},{\bf k}%
\right) -{\bf y}_{2}%
\end{array}%
\right)  \nonumber \\
&=&\left( 
\begin{array}{c}
{\Bbb A}_{12}{\bf x}_{2\#}\left( {\bf x}_{1},{\bf k}\right) -{\bf y}_{1} \\ 
0%
\end{array}%
\right)
\end{eqnarray}%
where 
\begin{eqnarray}
{\Bbb A}_{12} &=&{\Bbb A}_{12}\left( {\bf x}_{1},{\bf x}_{2\#}\left( {\bf x}%
_{1},{\bf k}\right) \right)  \nonumber \\
{\bf y}_{1} &=&{\bf y}_{1}\left( {\bf x}_{1},{\bf x}_{2\#}\left( {\bf x}_{1},%
{\bf k}\right) \right)
\end{eqnarray}%
Then solve the nonlinear equation for ${\bf x}_{1}$ 
\begin{eqnarray}
&&{\Bbb A}_{12}{\bf x}_{2\#}\left( {\bf x}_{1},{\bf k}\right) -{\bf y}_{1}=0
\nonumber \\
&\equiv &{\Bbb A}_{12}{\Bbb A}_{22}^{-1}\left( {\bf y}_{2}-{\Bbb A}_{21}{\bf %
x}_{1}\right) +{\Bbb A}_{12}{\bf k}-{\bf y}_{1}=0  \label{X1eq}
\end{eqnarray}%
with formal solution 
\[
{\bf x}_{1\#}={\Bbb A}_{21}^{-1}{\bf y}_{2}-{\Bbb A}_{21}^{-1}{\Bbb A}_{22}%
{\Bbb A}_{12}^{-1}{\bf y}_{1}={\Bbb A}_{21}^{-1}\left( {\bf y}_{2}-{\Bbb A}%
_{22}{\Bbb A}_{12}^{-1}{\bf y}_{1}\right) 
\]%
(which corresponds to the equation we identified for ${\cal V}_{den}$ in eq. 
$\left( \ref{gammaeq}\right) $) to determine the exact ${\bf x}_{1}$, where
everything in eq. $\left( \ref{X1eq}\right) $ is evaluated along the paths $%
\left( {\bf x}_{1},{\bf x}_{2\#}\left( {\bf x}_{1},{\bf k}\right) \right) $.
The different paths of solutions ${\bf x}_{2\#}\left( {\bf x}_{1},{\bf k}%
\right) $ are labelled by ${\bf k}$ i.e. 
\begin{equation}
{\bf x}_{2\#}\left( {\bf x}_{1},{\bf k}\right) \neq {\bf x}_{2\#}\left( {\bf %
x}_{1},{\bf k}^{\prime }\right) \text{ for }{\bf k}\left( {\bf x}_{1}\right)
\neq {\bf k}^{\prime }\left( {\bf x}_{1}\right)
\end{equation}%
and an extra criterion must be used to determine unique paths. As ${\bf x}%
_{2\#}\left( {\bf x}_{1},{\bf k}\right) $ is always a self consistent
solution for ${\bf x}_{2}$ in eq. $\left( \ref{x2eq}\right) $ one can
observe that for any ${\bf x}_{1}$ we have 
\begin{eqnarray}
&&\left( 
\begin{array}{cc}
0 & {\Bbb A}_{12} \\ 
{\Bbb A}_{21} & {\Bbb A}_{22}%
\end{array}%
\right) \left( 
\begin{array}{c}
{\bf x}_{1} \\ 
{\bf x}_{2\#}\left( {\bf x}_{1},{\bf k}\right)%
\end{array}%
\right) -\left( 
\begin{array}{c}
{\bf y}_{1} \\ 
{\bf y}_{2}%
\end{array}%
\right) =\left( 
\begin{array}{c}
{\Bbb A}_{12}{\bf x}_{2\#}\left( {\bf x}_{1}\right) -{\bf y}_{1} \\ 
{\Bbb A}_{21}{\bf x}_{1}+{\Bbb A}_{22}{\bf x}_{2\#}\left( {\bf x}_{1},{\bf k}%
\right) -{\bf y}_{2}%
\end{array}%
\right)  \nonumber \\
&=&\left( 
\begin{array}{c}
{\Bbb A}_{12}{\bf x}_{2\#}\left( {\bf x}_{1}\right) -{\bf y}_{1} \\ 
{\Bbb A}_{21}{\bf x}_{1}+{\Bbb A}_{22}\left( {\Bbb A}_{22}^{-1}\left( {\bf y}%
_{2}-{\Bbb A}_{21}{\bf x}_{1}\right) +{\bf k}\right) -{\bf y}_{2}%
\end{array}%
\right) =\left( 
\begin{array}{c}
{\Bbb A}_{12}{\bf x}_{2\#}\left( {\bf x}_{1}\right) -{\bf y}_{1} \\ 
{\bf 0}%
\end{array}%
\right)
\end{eqnarray}%
i.e.%
\begin{equation}
\left( 
\begin{array}{cc}
0 & {\Bbb A}_{12} \\ 
{\Bbb A}_{21} & {\Bbb A}_{22}%
\end{array}%
\right) \left( 
\begin{array}{c}
{\bf x}_{1} \\ 
{\bf x}_{2\#}\left( {\bf x}_{1}\right)%
\end{array}%
\right) -\left( 
\begin{array}{c}
{\bf y}_{1} \\ 
{\bf y}_{2}%
\end{array}%
\right) =\left( 
\begin{array}{c}
{\Bbb A}_{12}{\bf x}_{2\#}\left( {\bf x}_{1}\right) -{\bf y}_{1} \\ 
0%
\end{array}%
\right)
\end{equation}
